\scanpage{39}

D'où on conclut, utilisant l'équivalence de la catégorie
des rev. étales de $Y$ et des rev. étales de $X_{/Y}$, le

\uline{Corollaire} Si $(LG)$ vaut, alors le foncteur
$X'\longmapsto X'\times_XY$ des rev. étales de $X$ dans les
rev. étales de $Y$ est pleinement fidèle,
\deleted{Si $\unsure{X}$ $(LG)$ stricte vaut, pour}
et kif kif en remplaçant $\unsure{Y|X}$ par un voisinage
ouvert $U$ quelconque de $Y$.
Si $(LG)$ stricte vaut, alors pour tout rev étale $Y'$ de $Y$,
il existe un \deleted{rev étale $U'$ d'un}
voisinage ouvert $U$ de $Y$, et un rev. étale $U'$ de $U$,
\deleted{tel que} \inserted{\unsure{et d'un|et un}}
$Y$-isom $U'\times_UY\simeq Y'$.
(Si $U$ est choisi, $U'$ est bien déterminé à isom unique près).

\uline{Remarque}. Nous verrons plus loin des conditions,
liées au ``th de pureté'' permettant de conclure que
$\pi_1(U)\longrightarrow\pi_1(X)$ est
\deleted{surjectif}\inserted{bijectif}.
Kif-Kif pour Picard ---

\margin{38}
